\section{Mathematical Foundations}
\label{sec:method}

% Describe the methods and algorithms
% You need to explain how you implemented the methods and also say something
% about the structure of your algorithm and present some parts of your code. You
% should plug in some calculations to demonstrate your code, such as selected
% runs used to validate and verify your results. The latter is extremely
% important!! A reader needs to understand that your code reproduces selected
% benchmarks and reproduces previous results, either numerical and/or well-known
% closed form expressions.

% Should I structure this into three parts? The first being introducing
% mathematical methods and algorithms; the second, (high-level) implementation
% details for the code; and lastly, verifications to validate and verify the
% implementations results. Where to introduce the data set?

This section presents the mathematical methods used in this study. We first
introduce linear regression via ordinary least squares, followed by
regularization through Ridge and Lasso. We then discuss the bias--variance
trade-off, followed by the resampling methods bootstrap and $k$-fold
cross-validation, and conclude with gradient descent and automatic
differentiation. We assume familiarity with basic  linear algebra and
statistics.

\subsection{Linear regression}

Linear regression is among the earliest techniques in machine learning. It
addresses the scenario in which we are given $n$ samples with $p$ features,
arranged in a design matrix $\bm{X}\in \mathbb{R}^{n\times p}$, along with $n$
targets $\bm{y}\in \mathbb{R}^n$. We assume the targets are produced according
to
\begin{equation}\label{eq:data_generating_process}
    \bm{y} = f(\bm{x}) + \bm{\epsilon}, \qquad \bm{\epsilon}
    \sim \mathcal{N}(0, \bm{\sigma}^2),
\end{equation}
where $f$ denotes an unknown function, $\bm{\epsilon}$ is the noise vector, and
$\bm{\sigma}^2$ is the noise variance. The goal is to use the observed data to
build an estimator $\bm{\hat{y}}$ that approximates $f$.

We can define a design matrix as $\bm{X} \in \mathbb{R}^{n \times p}$, where $n$
is the number of instances and $p$ the number of parameters. We then also have
$\bm{\theta} \in \mathbb{R}^p$, which is the vector of model parameters
(weights). The predicted target values can then be expressed as
$\tilde{\bm{y}} = \bm{X}\bm{\theta}$, where $\tilde{\bm{y}} \in \mathbb{R}^n$.

\subsubsection{Singular value decomposition}
The singular value decomposition (SVD) is a powerful algorithm that allows us to
decompose any matrix $\bm{X}$ into a diagonal matrix $\bm{\Sigma}$ and two
orthogonal matrices, $\bm{U}$ and $\bm{V}$. The SVD can be expressed as
\begin{equation}
    \bm{X} = \bm{U\Sigma V^{\top}}.
\end{equation}
For a design matrix $\bm{X}$ of dimension $n \times p$, the matrix $\bm{U}$ is
$n \times n$, $\bm{V}$ is $p \times p$, and $\bm{\Sigma}$ is $n \times p$ with
$r = \min(n, p)$ singular values $\sigma_i \geq 0$ on the main diagonal and
zeros elsewhere~\cite{hjorth-jensenMachineLearningPhysical2026}.

\subsubsection{Ordinary least squares}

In ordinary least squares (OLS), the model parameters $\bm{\theta}$ are chosen
such that the predicted values $\bm{X}\bm{\theta}$ are as close as possible to
the observed response vector
$\bm{y} \in \mathbb{R}^n$~\cite{mehtaHighbiasLowvarianceIntroduction2019}.
This is achieved by minimizing the mean squared error (MSE),
\begin{equation}\label{eq:ols_cost}
    C_\text{OLS}(\bm{\theta}) = \frac{1}{n}\left\lVert \bm{X}\bm{\theta}
    - \bm{y} \right\rVert_2^2,
\end{equation}
where the factor $1/n$ does not affect the minimizer, but makes the objective
independent of sample size. Formally, we can denote the solution to this problem
as
\begin{equation}
    \bm{\hat{\theta}}_\text{OLS} = \underset{\bm{\theta}
    \in \mathbb{R}^p}{\arg\min}
    \left\lVert \bm{X}\bm{\theta} - \bm{y} \right\rVert_2^2,
\end{equation}
which after straightforward differentiation leads to the closed form equation
\begin{equation}\label{eq:ols-theta-hat}
    \bm{\hat{\theta}}_\text{OLS} = (\bm{X}^\top \bm{X})^{-1} \bm{X}^\top \bm{y}.
\end{equation}
The objective is convex, since its Hessian $\frac{2}{n}\bm{X}^\top \bm{X}$ is
positive semi-definite; any stationary point is therefore a global
minimizer~\cite{hjorth-jensenMachineLearningPhysical2026}.

The OLS estimator can also be expressed as
\begin{equation}
    \bm{\hat{\theta}}_\text{OLS} = \bm{X}^{+}\bm{y},
\end{equation}
where $\bm{X}^{+}$ denotes the Moore--Penrose pseudoinverse, which can be
computed using SVD \cite{goodfellowDeepLearning2016}. This formulation is
numerically more stable than solving the normal equations, since forming
$\bm{X}^{\top}\bm{X}$ squares the condition number of $\bm{X}$, which can result
in a substantial loss of numerical
precision~\cite{hjorth-jensenMachineLearningPhysical2026}.

% Check out Goodfellow et al (2016) Chapter 5.2.2 and Chapter 7, which covers
% regularization. This may be relevant for introducting Ridge and Lasso
% regression.

\subsubsection{Ridge regression}

By adding a regularizer defined as the $L_2$ norm of the the parameter vector to
the cost function, we achieve the penalized regression problem called Ridge
regression~\cite{mehtaHighbiasLowvarianceIntroduction2019}:
\begin{equation}
    \bm{\hat{\theta}}_{\text{Ridge}}(\lambda) = \underset{\bm{\theta}
    \in \mathbb{R}^p}{\arg\min}
    \left( \left\lVert \bm{X}\bm{\theta} - \bm{y} \right\rVert_2^2
    + \lambda \left\lVert \bm{\theta} \right\rVert_2^2 \right).
\end{equation}
Differentiating with respect to $\bm{\theta}$ yields a matrix inversion problem
similar to that of \cref{eq:ols-theta-hat}, but with a crucial difference: for
any finite $\lambda$, the added term removes the singularity issues that OLS can
encounter. The optimal parameters are then given by
\begin{equation}\label{eq:ridge-theta-hat}
    \bm{\hat{\theta}}_{\text{Ridge}}(\lambda) = (\bm{X^\top X}
    + \lambda \bm{I})^{-1} \bm{X^\top y}.
\end{equation}
This objective is strictly convex, meaning the minimum is unique, as the Hessian
$\bm{X^{\top}X}+\lambda\bm{I}$ is positive definite for all
$\lambda > 0$~\cite{hjorth-jensenMachineLearningPhysical2026}.

Just as with OLS, the Ridge prediction can be written using the SVD,
\begin{equation}
    \bm{\tilde{y}}_{\text{Ridge}} = \left( \sum_{j=0}^{p-1} d_j\,
    \bm{u}_j \bm{u}_j^{\top} \right) \bm{y},
\end{equation}
where each direction $\bm{u}_j$ is scaled by the shrinkage factor
\begin{equation}\label{eq:ridge_scaling}
    d_j = \frac{\sigma_j^2}{\sigma_j^2 + n\lambda} \in (0, 1].
\end{equation}
If $n\lambda \gg \sigma_j^2$ for some direction $j$, then
$d_j \approx \sigma_j^2/(n\lambda) \approx 0$, and that direction is almost
completely suppressed. Conversely, if $n\lambda \ll \sigma_j^2$, the factor
$d_j$ in \cref{eq:ridge_scaling} is just below one, and the direction is
essentially retained as in OLS.

% The Ridge estimator scales the OLS estimate by the inverse factor of $1+\lambda$,
% and the Ridge estimator converges to zero when $\lambda$ goes to infinity
% (\textbf{\textit{at least -- and probably only -- when the design matrix is orthonormal
%  -- I should look more into this}}).

\subsubsection{Lasso regression}

If we instead choose the $L_1$ norm for the regularizer, we obtain Lasso
regression, defined as~\cite{hjorth-jensenMachineLearningPhysical2026}
\begin{equation}\label{eq:lasso-theta-hat}
    \bm{\hat{\theta}}_{\text{Lasso}}(\lambda) = \underset{\bm{\theta}
    \in \mathbb{R}^p}{\arg\min} \left( \left\lVert \bm{X}\bm{\theta}
    - \bm{y} \right\rVert_2^2
    + \lambda \left\lVert \bm{\theta} \right\rVert_1 \right).
\end{equation}
Lasso stands for least absolute shrinkage and selection operator. Lasso has no
closed form solution, and can thus not be solved by a matrix inversion such as
\cref{eq:ols-theta-hat} and
\cref{eq:ridge-theta-hat}~\cite{hjorth-jensenMachineLearningPhysical2026}.
% Still a convex problem.

\subsection{Bias--variance decomposition}
\label{subsec:bias-variance}
% full derivation (required here, part c)

To decompose the bias--variance trade-off, we assume the data-generating process
of \cref{eq:data_generating_process}. The prediction $\tilde{\bm{y}}$ is a
random quantity because the parameters are fitted on a training set that is
itself random~\cite{hjorth-jensenMachineLearningPhysical2026}. All expectations
below are taken jointly over the noise $\bm{\varepsilon}$ on the test targets
and over the draw of the training set (equivalently, in the bootstrap, over the
resampled training sets).

Starting from the cost we wish to decompose (MSE) and adding and subtracting
$\mathbb{E}[\tilde{\bm{y}}]$ inside the square,
\begin{align}
    \mathbb{E}\!\left[(\bm{y}-\tilde{\bm{y}})^2\right]
    &=\mathbb{E}\!\left[\left(\bm{f}+\bm{\varepsilon}
                    -\tilde{\bm{y}}\right)^2\right]\nonumber\\
    &=\mathbb{E}\!\left[\left(
        (\bm{f}-\mathbb{E}[\tilde{\bm{y}}])
        +(\mathbb{E}[\tilde{\bm{y}}]-\tilde{\bm{y}})
        +\bm{\varepsilon}
    \right)^2\right].
\end{align}
Writing $a=\bm{f}-\mathbb{E}[\tilde{\bm{y}}]$,
$b=\mathbb{E}[\tilde{\bm{y}}]-\tilde{\bm{y}}$ and $c=\bm{\varepsilon}$, the
square expands into three squared terms and three cross terms,
\begin{align}
    \mathbb{E}\!\left[(a+b+c)^2\right]
    &= \mathbb{E}[a^2]+\mathbb{E}[b^2]+\mathbb{E}[c^2] \notag\\
    &\quad + 2\,\mathbb{E}[ab]+2\,\mathbb{E}[ac]+2\,\mathbb{E}[bc].
\end{align}
Consider the three cross terms:
\begin{align}
    \mathbb{E}[ab]
    &=\left(\bm{f}-\mathbb{E}[\tilde{\bm{y}}]\right)
    \mathbb{E}\!\left[\mathbb{E}[\tilde{\bm{y}}]-\tilde{\bm{y}}\right]=0,\\
    \mathbb{E}[ac]
    &=\left(\bm{f}-\mathbb{E}[\tilde{\bm{y}}]\right)\,\mathbb{E}[\bm{\varepsilon}]=0,\\
    \mathbb{E}[bc]
    &=\mathbb{E}\!\left[\mathbb{E}[\tilde{\bm{y}}]-\tilde{\bm{y}}\right]\mathbb{E}[\bm{\varepsilon}]=0,
\end{align}
where we used that $a$ is deterministic, that $\mathbb{E}[b]=\mathbb{E}[\mathbb{E}[\tilde{\bm{y}}]-\tilde{\bm{y}}]=\bm{0}$, that $\mathbb{E}[\bm{\varepsilon}]=\bm{0}$, and that $\bm{\varepsilon}$ is independent of both $\bm{f}$ and the training data (so it factors in $\mathbb{E}[bc]$). What remains is
\begin{equation}\label{eq:bias-variance-decomposition}
    \mathbb{E}\!\left[(\bm{y}-\tilde{\bm{y}})^2\right]
    =\underbrace{\left(\bm{f}-\mathbb{E}[\tilde{\bm{y}}]\right)^2}_{\mathrm{Bias}^2[\tilde{\bm{y}}]}
    +\underbrace{\mathbb{E}\!\left[\left(\tilde{\bm{y}}-\mathbb{E}[\tilde{\bm{y}}]\right)^2\right]}_{\mathrm{Var}[\tilde{\bm{y}}]}
    +\sigma^2 .
\end{equation}
In practice we report the sample averages over $n$ test points,
\begin{align}
    \mathrm{Bias}^2[\tilde{\bm{y}}]
    &= \frac{1}{n}\sum_i\left(f_i
    - \mathbb{E}[\tilde{y}_i]\right)^2, \label{eq:bias-term}\\
    \mathrm{Var}[\tilde{\bm{y}}]
    &= \frac{1}{n}\sum_i\left(\tilde{y}_i
    - \mathbb{E}[\tilde{y}_i]\right)^2, \label{eq:var-term}
\end{align}
with $\mathbb{E}[\tilde{y}_i]$ estimated by averaging the prediction at point
$i$ over the bootstrap replicates.

The three terms have distinct meanings. The \emph{bias} measures how far the
average model prediction is from the true function $f$: it is large when the
model is too rigid to represent $f$ (low polynomial degree). The \emph{variance}
measures how much the prediction fluctuates around its own mean as the training
set changes: it grows with model complexity. The \emph{noise variance}
$\sigma^2$ is irreducible and sets a floor on the achievable test
error~\cite{hjorth-jensenMachineLearningPhysical2026}.

Since $f$ is unknown, we cannot evaluate the bias directly and replace $\bm{f}$
by the observed targets $\bm{y}=\bm{f}+\bm{\varepsilon}$. The measured bias then
becomes
\begin{equation}\label{eq:measured-bias}
    \frac{1}{n}\sum_i\mathbb{E}[(y_i-\mathbb{E}[\tilde{y}_i])^2]
    =\mathrm{Bias}^2[\tilde{\bm{y}}]+\sigma^2,
\end{equation}
so the estimated bias absorbs the noise variance: the measured
$\mathrm{Bias}^2+\mathrm{Var}$ sums to the full test MSE, with $\sigma^2$ folded
into the bias rather than appearing as a separate
term~\cite{hjorth-jensenMachineLearningPhysical2026}.

\subsection{Resampling methods}

The standard situation in machine learning research is one where you have
limited access to data -- as opposed to our situation where we can easily
generate more as we know $f$.

Resampling methods are statistical techniques that repeatedly draw new random
samples from an existing dataset.

\subsubsection{Bootstrap}

\subsubsection{k-fold cross-validation}

% Goodfellow et al. (2016) Chapter 5.3.1 covers cross-validation

In $k$-fold cross-validation, the dataset is partitioned into $k$ folds, and the
model is trained $k$ separate times, each time leaving out a different fold for
evaluation~\cite{geronHandsonMachineLearning2023}.

This is commonly used in order to select hyperparameters, such as the penalty
parameter $\lambda$ in Ridge regression, in which you choose the value that
minimize the average estimate of the prediction error over the
folds~\cite{hjorth-jensenMachineLearningPhysical2026}.

\subsection{Automatic differentiation and backpropagation}
% Where do we discuss/introduce backpropagation and automatic differentiation?

\subsection{Gradient descent}

% The stability bound is set by lambda max and the convergence rate by lambda
% min. When they differ by orders of magnitude, the step that is barely safe
% along the steep direction is hopelessly small along the flat one. Steep
% coordinates need small steps; flat coordinates could take large ones.

%\[
%    \kappa(\bm{H}) = \kappa_2(\bm{X})^2
%\]

For optimization problems where an analytical (closed-form) solution is
unavailable, numerical optimization methods can be used to iteratively approach
a minimum of the cost function. One such method is Newton's method, which uses
both the gradient and the Hessian:
\begin{equation}\label{eq:newtons_method}
    \bm{\theta}^{(k+1)}
    =
    \bm{\theta}^{(k)}
    -
    \bm{H}^{-1}
    \nabla_{\bm{\theta}}
    C\!\left(\bm{\theta}^{(k)}\right).
\end{equation}
However, Newton's method requires computing and manipulating the Hessian. For a
model with $p$ parameters, the Hessian contains $p^2$ entries and inverting it
costs $\mathcal{O}(p^3)$, making it impractical to store and manipulate for
large-scale models~\cite{hjorth-jensenMachineLearningPhysical2026}.

Gradient descent provides a computationally simpler alternative by updating the
parameters using only the gradient: % Newton's method costs
% $\mathcal{O}(p^3)$~\cite{hjorth-jensenMachineLearningPhysical2026}
\begin{equation}\label{eq:gd_basic}
    \bm{\theta}^{(k+1)}
    =
    \bm{\theta}^{(k)}
    -
    \gamma
    \nabla_{\bm{\theta}}
    C\!\left(\bm{\theta}^{(k)}\right),
\end{equation}
where $\gamma>0$ is the learning rate which replaces the Hessian. This makes
gradient descent particularly suitable for large-scale machine-learning models,
such as neural networks, where the cost function is generally nonlinear and does
not have a closed-form solution.

\subsubsection{Momentum}

The convergence of gradient descent can be improved by adding a momentum term,
which carries over part of the previous parameter update. The resulting update
rule is
\begin{equation}\label{eq:gd_momentum}
    \bm{\theta}^{(k+1)} = \bm{\theta}^{(k)} -
    \gamma \nabla_{\bm{\theta}} C\!\left(\bm{\theta}^{(k)}\right)
    + \delta\!\left[\bm{\theta}^{(k)} - \bm{\theta}^{(k-1)}\right],
\end{equation}
where $\delta \in [0, 1]$ is the momentum parameter. This is known for
accelerating learning, especially in of high curvature, small but consistent
gradients, or noisy gradients. Momentum is mainly intended to address two
issues: ill-conditioning of the Hessian matrix and variance in the stochastic
gradient~\cite{goodfellowDeepLearning2016}. % Momentum amplifies the consistent
% and cancels the oscillatory~\cite{hjorth-jensenMachineLearningPhysical2026}.
% ``The method accelerates precisely where plain gradient descent
% crawls.~\cite{hjorth-jensenMachineLearningPhysical2026}''

\subsubsection{Stochastic gradient descent}
% Goodfellow et al. (2016) Chapter 5.9 and Hjorth-Jensen Chapter 4.7

For a cost $C(\bm{\theta}) = \frac{1}{n}\sum_{i=1}^{n} c_i(\bm{\theta})$, the
full gradient requires a pass over all $n$ data points per iteration, which is
expensive for large $n$~\cite{hjorth-jensenMachineLearningPhysical2026}.
Stochastic gradient descent (SGD) instead uses a mini-batch $\mathcal{B}_k$ of
$m$ examples i.i.d. sampled from the data, giving the unbiased gradient
estimate~\cite{goodfellowDeepLearning2016}
\[
    \bm{g}^{(k)} = \frac{1}{m} \sum_{i \in \mathcal{B}_k}
    \nabla_{\bm{\theta}} c_i(\bm{\theta}^{(k)}),
    \qquad
    \bm{\theta}^{(k+1)} = \bm{\theta}^{(k)} - \gamma \bm{g}^{(k)},
\]
where $\bm{g}^{(k)}$ is the gradient of the cost function at iteration $k$.

The sampling noise prevents SGD from converging exactly to the minimum; instead,
the iterates fluctuate around it with a spread scaling as $\sqrt{\gamma / m}$.
This can be reduced by increasing $m$, at a higher cost per iteration, or by
gradually decreasing $\gamma$ during training, known as learning rate
scheduling~\cite{hjorth-jensenMachineLearningPhysical2026}.

\subsubsection{AdaGrad}

AdaGrad, introduced by \textcite{duchiAdaptiveSubgradientMethods}, adapts the
step size of each coordinate individually by accumulating the squared gradients,
\[
    \bm{r}^{(k)} = \bm{r}^{(k-1)} + \bm{g}^{(k)} \odot \bm{g}^{(k)},
    \qquad \bm{r}^{(0)} = \bm{0},
\]
where $\odot$ denotes the Hadamard product. The parameters are then updated as
\[
    \bm{\theta}^{(k+1)} = \bm{\theta}^{(k)}
    - \frac{\gamma}{\sqrt{\bm{r}^{(k)}} + \epsilon} \odot \bm{g}^{(k)},
\]
with the square root and division taken element-wise, and $\epsilon$ a small
constant for numerical stability. Coordinate $j$ thus has the effective step
size given by
\[
    \frac{\gamma} {\sqrt{r_{j}^{(k)}} + \epsilon},
\]
which decreases monotonically as gradients
accumulate~\cite{hjorth-jensenMachineLearningPhysical2026}. Coordinates with
large past gradients therefore take smaller steps, while rarely updated
coordinates take larger steps~\cite{goodfellowDeepLearning2016}.

\subsubsection{RMSProp}
% Introduced in Tieleman & Hinton (2012), Coursera Lecture 6.5;
% see also Goodfellow et al. (2016) Sec. 8.5.2

Since AdaGrad accumulates all past squared gradients, its effective step size
decreases monotonically, which can halt learning prematurely. RMSProp addresses
this by replacing the cumulative sum with an exponentially decaying
average~\cite{goodfellowDeepLearning2016},
\[
    \bm{v}^{(k)} = \rho\, \bm{v}^{(k-1)} + (1-\rho)\,
    \bm{g}^{(k)} \odot \bm{g}^{(k)},
    \qquad \bm{v}^{(0)} = \bm{0},
\]
where $\rho \in [0, 1)$ is the decay rate, typically $0.9$ or $0.99$, giving an
effective memory of about $1/(1-\rho)$ iterations. The parameters are updated as
\[
    \bm{\theta}^{(k+1)} = \bm{\theta}^{(k)}
    - \frac{\gamma}{\sqrt{\bm{v}^{(k)}} + \epsilon} \odot \bm{g}^{(k)},
\]
with the square root and division taken
element-wise~\cite{hjorth-jensenMachineLearningPhysical2026}.

\subsubsection{Adam}

\cite{kingmaAdamMethodStochastic2017}

% Learning rate decay?
